\documentclass[11pt,a4paper]{article}
\newcommand{\fichetype}{Exercices}
% Thème A4 commun, librement modifiable.
\usepackage[T1]{fontenc}
\usepackage[utf8]{inputenc}
\usepackage[french]{babel}
\usepackage{lmodern,amsmath,amssymb,mathtools,tikz,pgfplots}
\usepackage[margin=2.2cm,headheight=15pt]{geometry}
\usepackage{xcolor,fancyhdr,enumitem,hyperref,needspace}
\usetikzlibrary{arrows.meta,positioning,calc}
\pgfplotsset{compat=1.18,width=11cm,height=6cm}
\definecolor{Ink}{HTML}{202124}
\definecolor{Blue}{HTML}{356B9E}
\definecolor{Muted}{HTML}{737579}
\color{Ink}
\renewcommand{\familydefault}{\sfdefault}
\setlength{\parindent}{0pt}
\setlength{\parskip}{6pt}
\setlist{itemsep=3pt,topsep=4pt}
\pagestyle{fancy}\fancyhf{}
\fancyhead[L]{\small\color{Muted}MATHS LYCÉE}
\fancyhead[R]{\small\color{Muted}2026--2027}
\fancyfoot[L]{\small\color{Muted}\fichetype}
\fancyfoot[R]{\small\thepage}
\renewcommand{\headrulewidth}{0.3pt}
\hypersetup{colorlinks=true,urlcolor=Blue,linkcolor=Blue}
\newcommand{\R}{\mathbb{R}}
\newcommand{\N}{\mathbb{N}}
\newcommand{\E}{\mathbb{E}}
\newcommand{\V}{\mathrm{V}}
\newcommand{\exercice}[3]{\Needspace{6\baselineskip}\vspace{0.5cm}{\color{Blue}\small\textbf{EXERCICE #1}\hfill #2}\par{\large\bfseries #3}\par}

\begin{document}
{\small\color{Blue}PREMIÈRE SPÉCIALITÉ}\par{\LARGE\bfseries Produit scalaire}\par{\large Exercices}\par\vspace{0.2cm}\hrule\vspace{0.4cm}
Les exercices progressent de l’application directe vers la recherche. Détailler les calculs et justifier les réponses.
\par\noindent\begin{minipage}{\linewidth}\exercice{1}{Application}{Calculs de base}
Dans un repère orthonormé, $\vec u=(3,-1)$, $\vec v=(2,4)$ et $\vec w=(-1,2)$. Calculer $\vec u\cdot\vec v$, $\|\vec u\|$, puis $\vec u\cdot(2\vec v-\vec w)$ de deux façons.
\end{minipage}\par
\par\noindent\begin{minipage}{\linewidth}\exercice{2}{Application}{Projection et signe}
$OA=6$. Le projeté orthogonal $H$ de $B$ sur $(OA)$ vérifie $OH=4$. Calculer $\overrightarrow{OA}\cdot\overrightarrow{OB}$ lorsque $H\in[OA)$, puis lorsque $H$ est sur la demi-droite opposée. Que peut-on dire du type de l’angle $\widehat{AOB}$ ?
\end{minipage}\par
\par\noindent\begin{minipage}{\linewidth}\exercice{3}{Application}{Un triangle rectangle}
On donne $A(1,2)$, $B(4,3)$ et $C(0,5)$. Démontrer que le triangle est rectangle. Calculer ses trois longueurs et son aire.
\end{minipage}\par
\par\noindent\begin{minipage}{\linewidth}\exercice{4}{Entraînement}{Paramètre et orthogonalité}
Pour $\vec u=(2,t)$ et $\vec v=(t-1,3)$, déterminer $t$ pour que les vecteurs soient orthogonaux. Pour cette valeur, calculer leurs normes.
\end{minipage}\par
\par\noindent\begin{minipage}{\linewidth}\exercice{5}{Entraînement}{Deux côtés et un angle}
Dans un triangle, $AB=4$, $AC=7$ et $\widehat{BAC}=120^\circ$. Déterminer exactement $BC$. Puis calculer $\overrightarrow{AB}\cdot\overrightarrow{AC}$.
\end{minipage}\par
\par\noindent\begin{minipage}{\linewidth}\exercice{6}{Entraînement}{Retrouver un angle}
Un triangle vérifie $AB=5$, $AC=6$ et $BC=7$. Calculer $\cos\widehat{BAC}$, puis donner l’angle au degré près. Est-il rectangle ?
\end{minipage}\par
\par\noindent\begin{minipage}{\linewidth}\exercice{7}{Approfondissement}{Un lieu géométrique}
Soient $A(-2,1)$ et $B(4,1)$. Déterminer l’ensemble des $M(x,y)$ vérifiant $\overrightarrow{MA}\cdot\overrightarrow{MB}=0$. Donner centre et rayon, puis tester $C(1,4)$.
\end{minipage}\par
\par\noindent\begin{minipage}{\linewidth}\exercice{8}{Approfondissement}{Automatiser un test géométrique}
On veut tester avec un programme si trois points distincts $A,B,C$ forment un angle droit en $A$. Décrire les calculs nécessaires. Appliquer à $A(0,0)$, $B(6,2)$, $C(-1,3)$. Pourquoi comparer exactement à zéro peut-il poser problème avec des décimaux approchés ?
\end{minipage}\par
\end{document}
